\documentclass[11pt,a4paper]{article}
\usepackage[T1]{fontenc}
\usepackage{lmodern,amsmath,amssymb,amsthm,mathtools,booktabs,microtype}
\usepackage[margin=27mm]{geometry}
\usepackage[colorlinks=true,linkcolor=blue,citecolor=blue,urlcolor=blue]{hyperref}
\setlength{\emergencystretch}{2em}
\newtheorem{theorem}{Theorem}[section]
\newtheorem{lemma}[theorem]{Lemma}
\newtheorem{proposition}[theorem]{Proposition}
\newtheorem{corollary}[theorem]{Corollary}
\theoremstyle{remark}\newtheorem{remark}[theorem]{Remark}
\newcommand{\Q}{\mathbb Q}\newcommand{\Qbar}{\overline{\mathbb Q}}
\DeclareMathOperator{\Tr}{Tr}\DeclareMathOperator{\Res}{Res}
\DeclareMathOperator{\divisor}{div}\DeclareMathOperator{\Pic}{Pic}
\title{Ramification of fixed-argument native complement covers}
\author{}\date{16 November 2026}
\begin{document}\maketitle
\begin{abstract}
Fix the input, root points, comparison scales, and slope of a native
detector, with positive scalar coefficient. Its square-complement
double cover of the full independent auxiliary-point variety is
geometrically ramified. It therefore has no rational section after
pullback by a dominant rational map from an abelian variety, including
the multiplication-by-two map used to represent doubled source points.
On a geometric boundary divisor, the trace detector restricts to the
reciprocal of an elliptic function with three poles, each of order four.
An unramified complement cover would restrict that function to a finite
constant set, whereas a native index-comparison factor has moving zeros.
The result does not exclude rational lifts at individual tuples or on
proper subvarieties of the auxiliary-point space.
\end{abstract}

\section{Fixed data and independent auxiliary points}

Let
\[
 K=\Q(\sqrt2),\qquad d=3-\sqrt2,\qquad
 E:\ y^2=x^3-d^2x,\qquad \mathcal E=\Res_{K/\Q}E,
\]
and let $\sigma$ be the nontrivial automorphism of $K$. Write
\[
 f=1/x,\qquad z=-x/y,
\]
and retain the fixed nonzero native point $P$.
Fix the rational input, all root points, the comparison scales, and
the comparison slope of a native source in \cite{UnitNorm}.
All resulting contact and scalar denominators are nonzero. The fixed
root point denoted $A\in E(K)$ has nonzero abscissa and ordinate,
as required by its native guards. In particular it is not a two-torsion
point. The other fixed root guards are retained as well.

For a general source the auxiliary roles are
$H,U_A,U_S,U_T,U_Z$, independently varying on
\[
                         \mathcal B=\mathcal E^5.
\]
For either conditional source they are $H,U$, independently varying
on $\mathcal B=\mathcal E^2$. We assume the native guard functions
on this variety are nonzero rational functions. This includes any
scalar denominator that becomes constant after the root data are fixed.

For one detector put
\begin{equation}\label{eq:incidence}
 V=\Tr_{K/\Q}(\mathcal W^{-1}),\qquad
 b=\frac{4w}{V^2},\qquad k=\frac{4w}{17},
 \qquad \mathcal R=V^2-k.
\end{equation}
Here $w$ is fixed and positive. In the general source $w=\rho R$;
in the conditional source $w=F/(4\delta^8)$. Positivity follows
whenever the earlier data admit a prepared native tuple, since the
input is $w$ times a nonzero square. The square-complement equation is
equivalent to the rational incidence equation $Y^2=\mathcal R$:
\[
             17(17-b)=\left(\frac{17}{V}\right)^2\mathcal R.
\]

\begin{theorem}\label{thm:main}
Under these fixed-data hypotheses, the function $\mathcal R$ has
odd valuation at some geometric prime divisor of $\mathcal B$.
Consequently its quadratic extension is geometrically ramified.

For every dominant rational map $\phi:B\dashrightarrow\mathcal B$
from an abelian variety $B/\Q$, the pullback $\phi^*\mathcal R$
is not a square in $\Q(B)$. In particular, the complement incidence
has no rational section after any isogeny, group translation followed
by a surjective homomorphism, or multiplication by two on all auxiliary
point factors.
\end{theorem}

The theorem concerns a section on the full independent auxiliary-point
space and its dominant abelian pullbacks. The detector input $b$ varies
with the auxiliary points; its square class is the fixed class of $w$.
The theorem does not
exclude rational lifts at selected tuples, a lift on a proper
subvariety, or a direct common-pole construction that does not force
square complements. In particular it does not resolve the existence
problem for the prepared independent-index families of \cite{Returns}.

\section{The native boundary formula}

\subsection{Dependence on one auxiliary role}

For the general source write $Q_1=A,Q_2=S,Q_3=T,Q_4=Z$,
$U_0=H$, $V_0=H-P$, and $V_i=U_i-Q_i$. Put
\[
 L_i=\frac{z(U_i)}{z(H)},\qquad
 L'_i=\sigma\left(\frac{z(V_i)}{z(H-P)}\right),\qquad
 \beta=216X^4+450459.
\]
The sixteen native field tests are
\begin{equation}\label{eq:generaltests}
 (f(U_i))_{i=0}^4,\quad (\sigma f(V_i))_{i=0}^4,\quad
 (L_i-L'_i)_{i=1}^4,\quad L_1-\beta,\quad
 \mathcal F(L_2,L_3,L_4).
\end{equation}
Their orders $a_i$ are one except at positions $12,13,14,16$,
which have order three. The six scalar tests $S_\nu$ are the two
rational components of $f(H)$, $f(S-P)$, and $z(S-P)/z(H)$.
Their orders are $\kappa_\nu=(1,1,1,1,3,3)$, and
\[
 G_\nu=R^{\kappa_\nu}\frac{RS_\nu}{1+RS_\nu}.
\]
In particular no scalar test contains $U_A$. With the fixed native
$\delta$ and $w$, the detector is
\begin{equation}\label{eq:generalW}
 \mathcal W=
 \prod_{i=1}^{16}\left(1+\frac{C_i^2}{\delta^2w^{2a_i-1}}\right)
 \prod_{\nu=1}^6\left(1+\frac{G_\nu^2}{w^{2\kappa_\nu-1}}\right).
\end{equation}

For the conditional source, put
\[
 L=\frac{z(U)}{z(H)},\qquad
 L'=\sigma\left(\frac{z(U-A)}{z(H-P)}\right).
\]
The six tests and product are
\begin{equation}\label{eq:conditionalW}
\begin{gathered}
 f(H),\ f(U),\ \sigma f(H-P),\ \sigma f(U-A),\ L-L',\ L-X,\\
 \mathcal W=\prod_{i=1}^{6}\left(1+\frac{\delta^4C_i^2}{F}\right).
\end{gathered}
\end{equation}
These are the direct index-comparison formulas of
\cite{Direct,UnitNorm}. Both weak-$abc$ parameter choices have the
same six tests.

\subsection{A geometric boundary divisor}

Over $\Qbar$ split $\mathcal E=E_+\times E_-$ and use subscripts
for the corresponding point coordinates and detector factors.
Let $T_0=(0,0)$ on $E_-$. Set
\[
                 U_{A,-}=T_0
\]
in the general source, or $U_-=T_0$ in the conditional source.
This defines a geometric prime divisor $D_0$ in $\mathcal B_{\Qbar}$.
It lies outside the guarded affine source open set; only the
restriction of rational functions at its generic point will be used.

For any point $Q=(x,y)$ with $xy\ne0$, the addition law gives
\begin{equation}\label{eq:torsiontranslate}
 x(T_0-Q)=-\frac{d^2}{x},\qquad
 y(T_0-Q)=-\frac{d^2y}{x^2},\qquad
 z(T_0-Q)=z(Q).
\end{equation}
The analogous formula holds with the conjugate value of $d$.
Thus the native nonzero abscissa and ordinate of $A$ ensure that
$T_0-A_-$ is finite and avoids every pole of $z$ and $f$ that
occurs in this restriction.

The factor containing $f(U_{A,-})^2$, or $f(U_-)^2$, has a pole
of order four along $D_0$. Every other factor is generically a unit
there. Indeed, after the displayed substitution its remaining
denominators can be avoided on a nonempty real open set of the
independent point roles. The fixed denominators are nonzero by
hypothesis, the difference involving $A$ is controlled by
\eqref{eq:torsiontranslate}, and the remaining variable denominators
are avoided by generic choices of $H$ and the other auxiliary roles.
At such a real tuple each remaining detector factor is one plus a
nonnegative quantity, since $w>0$. Hence it cannot vanish
identically on $D_0$. It follows that
\begin{equation}\label{eq:residue}
 v_{D_0}(\mathcal W_-)=-4,\qquad
 v_{D_0}(\mathcal W_+)=0,\qquad
                 V|_{D_0}=W^{-1},
\end{equation}
where $W=\mathcal W_+|_{D_0}$.

Write $U=U_{A,+}$ or $U=U_+$, respectively, and put
\[
 h=z(H_+),\qquad
 c=\frac{z(T_0-A_-)}{z(H_--P_-)}.
\]
Exactly three positive-embedding field tests still depend on $U$:
\[
                 f(U),\qquad z(U)/h-c,\qquad z(U)/h-a.
\]
Their positions and constants are
\begin{center}
\begin{tabular}{llll}\toprule
Source & Positions & $a$ & $\alpha$\\\midrule
General & $2,11,15$ & $\beta$ & $1/(\delta^2w)$\\
Conditional & $2,5,6$ & $X$ & $\delta^4/F=1/(4\delta^4w)$\\\bottomrule
\end{tabular}
\end{center}
In the general source all three positions have order one.
Both values of $\alpha$ are positive. Consequently
\begin{equation}\label{eq:Wboundary}
 W(U)=K_0\bigl(1+\alpha f(U)^2\bigr)
       \bigl(1+\alpha(z(U)/h-c)^2\bigr)
       \bigl(1+\alpha(z(U)/h-a)^2\bigr),
\end{equation}
where $K_0\ne0$ is independent of $U$. It may depend on $H$
and on all other auxiliary roles; it is not assumed constant in
those parameters.

\section{A fixed pole divisor and moving zeros}

\begin{lemma}\label{lem:divisor}
As a rational function on the fixed elliptic curve $E_+$, over
the function field of the other parameters on $D_0$, the function
$W$ has pole divisor
\begin{equation}\label{eq:poledivisor}
                   \Delta=4(T_0+T_++T_-),
\end{equation}
where $T_0=(0,0)$ and $T_\pm=(\pm d_+,0)$ on $E_+$.
The family $W$ is not a constant family of rational functions on
$E_+$.
\end{lemma}
\begin{proof}
The orders of the elementary functions are
\begin{center}
\begin{tabular}{lrrrr}\toprule
Point & $v(x)$ & $v(y)$ & $v(f)$ & $v(z)$\\\midrule
$O$ & $-2$ & $-3$ & $2$ & $1$\\
$T_0$ & $2$ & $1$ & $-2$ & $1$\\
$T_+,T_-$ & $0$ & $1$ & $0$ & $-1$\\\bottomrule
\end{tabular}
\end{center}
At $T_0$, only the first factor of \eqref{eq:Wboundary} has a
pole, of order four. Its remaining leading coefficient contains
\[
                  (1+\alpha c^2)(1+\alpha a^2),
\]
which is generically nonzero by real positivity. At each $T_\pm$,
the two factors involving $z$ each have a pole of order two,
while the first factor takes the nonzero value
$1+\alpha/d_+^2$. There are no other poles. This proves
\eqref{eq:poledivisor}.

The last factor in \eqref{eq:Wboundary} has zeros at
\begin{equation}\label{eq:moving}
                    z(U)=h(a\pm i/\sqrt\alpha).
\end{equation}
Both multipliers are nonzero, since $a$ is real and $\alpha>0$.
The function $h=z(H_+)$ is nonconstant, independently of $U$.
Thus these zeros move as $H_+$ varies. They cannot cancel against
$K_0$, which has no $U$ dependence, or against poles, whose
support is the fixed three-point set in \eqref{eq:poledivisor}.
The family of functions $W$ is therefore nonconstant.
\end{proof}

\section{Finiteness for an unramified complement cover}

\begin{lemma}\label{lem:finite}
Let $k_0$ be an algebraically closed field of characteristic zero,
let $E/k_0$ be an elliptic curve, and let $T_0,T_+,T_-$ be its
three distinct nonidentity two-torsion points. Fix
$M\in k_0^\times$ and $\Delta=4(T_0+T_++T_-)$. There are
only finitely many $W\in k_0(E)$ with pole divisor exactly
$\Delta$ for which $\divisor(W^2-M^2)$ is even.
After any algebraically closed extension of $k_0$, every such
function still belongs to this same finite set.
\end{lemma}
\begin{proof}
First suppose $W^2-M^2=R_0^2$ in the elliptic function field.
Set $Z=W+R_0$. Then
\[
                   W=\frac{Z+M^2/Z}{2}.
\]
The poles of $W$ are exactly the zeros and poles of $Z$, with
their orders unchanged. Their three equal weights four cannot
be partitioned into zero and pole divisors of equal degree.
Thus the square case is impossible, over any extension field.

If $\divisor(W^2-M^2)$ is even, adjoining its square root instead
gives a connected etale double cover $E'\to E$. There are only
three such covers over an algebraically closed field. Indeed their
trace-zero invertible sheaves are the three nontrivial elements of
$\Pic^0(E)[2]$, and their multiplication is unique up to scaling
by a square constant. All three covers and their points above the
$T_i$ are defined over $k_0$ and remain the only choices after
algebraically closed extension. Each $E'$ is an elliptic curve
for a suitable origin; the covering is an isogeny
\cite[Remark~15.3]{Milne}.

Let $\iota$ be its deck involution and put
\[
               Z=W+\sqrt{W^2-M^2},\qquad Z\iota(Z)=M^2.
\]
Above each pole of $W$ of order four, $Z$ has a zero of order
four at one preimage and a pole of order four at the other.
It has no other zeros or poles. There are consequently only
finitely many possible divisors of $Z$ on the three fixed curves.

For any such principal divisor choose a function $Z_0$ over $k_0$.
This choice is possible even if principality was first obtained
after extending the constants: on a constant elliptic curve, a
degree-zero divisor supported at constant points is principal
exactly when its weighted sum in the elliptic group is zero.
That is already an equality of points over $k_0$.
Every other function with the same divisor is $cZ_0$ for a
nonzero constant $c$. Since $Z_0\iota(Z_0)$ is a nonzero constant,
the norm condition gives
\[
                      c^2=\frac{M^2}{Z_0\iota(Z_0)}.
\]
There are at most two choices, both in $k_0$. Finally
$W=(Z+M^2/Z)/2$ is determined. This proves both finiteness and
its constant-field assertion.
\end{proof}

The special pole divisor is essential. For a general divisor whose
weights can be balanced, the square case can admit the varying family
$(cZ_0+M^2/(cZ_0))/2$. Lemma~\ref{lem:finite} excludes that case
by its three equal pole weights.

\section{Ramification and dominant abelian pullbacks}

\begin{proof}[Proof of Theorem~\ref{thm:main}]
Work over $\Qbar$. Suppose that every divisor valuation of
$\mathcal R$ on $\mathcal B_{\Qbar}$ were even. On this smooth
projective variety, write $\divisor(\mathcal R)=2D$. The
trivialization of $\mathcal O(2D)$ then defines a finite etale
double cover with function algebra obtained by adjoining
$\sqrt{\mathcal R}$; a split cover is allowed. Locally the
construction has equation $u^2=v$ with $v$ a unit, so there
is no missing ramification in higher codimension.

Restrict this etale cover to $D_0$. At its generic point,
\eqref{eq:residue} gives
\[
                    \mathcal R|_{D_0}=W^{-2}-k,
\]
a nonzero rational function because $W$ is nonconstant in $U$.
Next take the generic fiber of the projection of $D_0$ onto all
its elliptic factors except $U$. This fiber is the constant curve
$E_+$ over the function field of the remaining independent
parameters. The restricted double cover remains etale, including
after algebraic closure of that coefficient field.

Multiplication of its radicand by a square and by a nonzero
algebraic constant identifies its function extension with
\[
                         \sqrt{W^2-M^2},\qquad
                         M^2=\frac1k=\frac{17}{4w}.
\]
Lemma~\ref{lem:divisor} supplies the fixed pole divisor
$4(T_0+T_++T_-)$, so Lemma~\ref{lem:finite} would force $W$
to belong to a finite constant set of elliptic functions.
This contradicts its moving zero divisor. Some geometric divisor
valuation of $\mathcal R$ must therefore be odd.

Now let $\phi:B\dashrightarrow\mathcal B$ be a dominant
rational map from an abelian variety. It extends to a morphism
and is a translation of a homomorphism
\cite[Theorem~3.1 and Corollary~2.2]{Milne}. Dominance makes it
surjective, and it is smooth in characteristic zero. The inverse
image of a reduced prime divisor under a smooth surjective map
has multiplicity one at each of its generic points. An odd
valuation of $\mathcal R$ therefore remains odd along a component
of its inverse image. Thus $\phi^*\mathcal R$ is not a square.
\end{proof}

The native auxiliary points are obtained as doubles of independent
half-points. Multiplication by two is a finite etale surjective map
of $\mathcal E^5$ or $\mathcal E^2$. The ramification proved above
therefore persists on precisely that half-point parameter space;
the argument does not replace doubled source points by unconstrained
scalar coefficients.

The conclusion is geometric nonexistence of a section. A rational
point of a ramified cover can still project to a suitable native
tuple, and a proper coupled subvariety need not dominate
$\mathcal B$. These possibilities, as well as direct common-pole
constructions, remain outside the obstruction.

\section{Exact verification}

The reproducible standard-library checker is
\begin{center}\small
\texttt{negative\_results/computation/fixed\_argument\_complement\_v44/}\\
\texttt{check.py}.
\end{center}
It binds the native test lists and orders, verifies the torsion
translation and conditional scalar identities, audits the elliptic
orders and the leading coefficients at all three poles, and checks
the impossibility of balancing the three equal pole weights.
The constant-field finiteness argument, etale restriction, and
smooth pullback statements are proved above; the finite checks
do not assert a rational witness or a variable-count improvement.

\begin{thebibliography}{9}
\raggedright
\bibitem{Milne}
J. S. Milne, \emph{Abelian Varieties}, in
G. Cornell and J. H. Silverman (eds.), \emph{Arithmetic Geometry},
Springer, 1986, pp.~103--150.
Author's version dated 2 January 2022:
\url{https://www.jmilne.org/math/xnotes/AVs.pdf}.

\bibitem{UnitNorm}
\emph{A single norm determinant for the 25- and 13-variable constructions},
14 November 2026. Repository source:
\texttt{papers/finite\_algebra/unit\_norm\_20261114\_v28.tex}.

\bibitem{Direct}
\emph{Thirteen Variables from Direct Elliptic Index Comparison},
10 November 2026. Repository source:
\texttt{papers/conditional\_few\_variables/}\\
\texttt{direct\_input13\_20261110\_v22.tex}.

\bibitem{Returns}
\emph{Fixed pole fibers and independent detector indices},
16 November 2026. Repository source:
\texttt{papers/pole\_parameters/}\\
\texttt{common\_parameter\_guard\_reduction\_20261116\_v42.tex}.
\end{thebibliography}
\end{document}
